%& /Users/ziky/.cache/org-persist/0e/3b33a0-6345-4590-aeec-1fd24706a50c-1fc64c38fa07e6f8bfe4daf826445aa2
% Created 2025-09-08 Mon 10:34
% Intended LaTeX compiler: pdflatex
\documentclass[11pt]{article}
\usepackage[utf8]{inputenc}
\usepackage[T1]{fontenc}
\usepackage{amsmath}
\usepackage{amssymb}
\usepackage{capt-of}
\usepackage{hyperref}
\setlength{\abovedisplayskip}{0pt}
\setlength{\belowdisplayskip}{0pt}
\usepackage[a4paper, margin=1in]{geometry}

%% ox-latex features:
%   !announce-start, !guess-pollyglossia, !guess-babel, !guess-inputenc,
%   maths, !announce-end.

\usepackage{amsmath}
\usepackage{amssymb}

%% end ox-latex features


% end precompiled preamble
\ifcsname endofdump\endcsname\endofdump\fi

\author{Ziky Zhang}
\date{\today}
\title{Worksheet \#2}
\hypersetup{
 pdfauthor={Ziky Zhang},
 pdftitle={Worksheet \#2},
 pdfkeywords={},
 pdfsubject={},
 pdfcreator={},
 pdflang={English}}
\begin{document}

\maketitle
\newpage
Question set:
\begin{enumerate}
\item Two charges lie along the x-axis, a \(-8.0 \mu C\) charge at \(x=0\) and a \(-2.0 \mu C\) charge at \(x=6.0cm\).
\begin{enumerate}
\item Calculate the electric field at the point \((x, y)=(0, 8.0cm)\) (along the y-axis).
\item Find all points (besides infinity) where the electric field is zero.
\end{enumerate}
\item A charged rod is placed on the x-axis between \(x=d\) and \(x=d+L\). Its charge density is not uniform, and is given by \(\lambda=cx\), where \(c\) is a positive constant.
\begin{enumerate}
\item What are the SI units of \(c\)?
\item Calculate the total charge of the rod (hint: \(Q= \int dq\)).
\item Calculate the electric field at \(x=0\).
\item Suppose \(L=2d\). If the rod is replaced by a point charge with the same total charge (calculated in part b), where should the point charge be located in order to give the same electric field at \(x=0\) as for the rod.
\item A proton is placed at \(x=0\) and is released from rest. Describe the trajectory of the proton. Is the acceleration of the proton constant?
\end{enumerate}
\end{enumerate}

\begin{align*}
1. \ 
\text{Variables: }
q_{1} = -&8.0 \text{ at } (0,0) \\
q_{2} = -&2.0 \text{ at } (6,0) \\
        (&\mu C) \quad (cm) \\
\text{Vector: }
\hat{r_{1}} = &(0, 1) \\
\hat{r_{2}} = &(-\frac{6}{10}, \frac{8}{10})
\end{align*}

\begin{align*}
1.(a) \ 
\overrightarrow{E}(x,y) &= k_{e} \frac{q_{1}}{r_{1}^{2}} \hat{r_{1}} + k_{e} \frac{q_{2}}{r_{2}^{2}} \hat{r_{2}} \ \\
          &= k_{e} (\frac{q_{1}}{( \sqrt{0 + 8^{2}} cm )^{2}} \hat{r_{1}} + \frac{q_{{2}}}{( \sqrt{(-6)^{2} + 8^{2}} cm )^{2}} \hat{r_{2}}) \\
          &= k_{e} (\frac{-8.0 \mu C}{( 8 cm )^{2}} \hat{r_{1}} + \frac{-2.0 \mu C}{( 10 cm )^{2}} \hat{r_{2}}) \\
          &= - k_{e} (\frac{1}{8} \hat{r_{1}} + \frac{1}{50} \hat{r_{2}}) \frac{10^{-6} C}{10^{-4} m^{2}} \\
          &= - 8.99 \times 10^{7} \frac{N}{C} (\frac{1}{8} \hat{r_{1}} + \frac{1}{50} \hat{r_{2}}) \\
          &= - 1.12375 \times 10^{7} \hat{r_{1}} - 1.798 \times 10^{6} \hat{r_{2}} \frac{N}{C} \\
          &= ( (- 1.12375 \cdot 0)\hat{\imath} - 1.12375 \cdot 1 \hat{\jmath}) \times 10^{7} \\
& \quad + (- 1.798 \cdot -\frac{6}{10} \hat{\imath} - 1.798 \cdot \frac{8}{10} \hat{\jmath}) \times 10^{6} ) \frac{N}{C} \\
          &= ( (0 \hat{\imath} - 1.12375 \hat{\jmath}) \times 10^{7} + (1.0788 \hat{\imath} - 1.438\hat{\jmath})  \times 10^{6} ) \frac{N}{C} \\
          &= (1.08 \times 10^{6} \hat{\imath} - 1.268 \times 10^{7} \hat{\jmath}) \frac{N}{C}
\end{align*}

\newpage

1.(b) 
Since both charges in the system are laying on the \(x\)-axis, we know that the point with zero charge is also on the \(x\) -axis. So this part of this question can be simplified down to 1-D. \\
Say the location of the zero charge point of electric field as variable \(x\), therefore, we know the distance from \(q1\) and \(q2\) to the point \(p\) is \(x\) cm and \(6-x\) cm, respectively. \\
\begin{align*}
              E = 0 &= k_{e} \frac{| q_{1} |}{r_{1}^{2}} \hat{r_{1}} + k_{e} \frac{| q_{2} |}{r_{2}^{2}} \hat{r_{2}} \\
                    &= k_{e} ( \frac{| q_{1} |}{x^{2}} + \frac{| q_{2} |}{(6-x)^{2}} ) \\
                    &= \frac{q_{1}}{( x \text{ cm} )^{2}} + \frac{q_{2}}{( (6-x) \text{ cm} )^{2}} \\
- \frac{| q_{2} |}{( (6-x \text{ cm} )^{2}} &= \frac{| q_{1} |}{( x \text{ cm} )^{2}} \\
\bigg | -\frac{q_{2}}{q_{1}} \bigg | &= \frac{( (6-x) \text{ cm} )^{2}}{( x \text{ cm} )^{2}} \\
\sqrt{ \tfrac{q_{2}}{q_{1}} } &= \frac{6-x}{x} \text{ cm} \\
\sqrt{ \tfrac{q_{2}}{q_{1}} } &= \frac{6}{x} - 1 \\
\sqrt{ \tfrac{q_{2}}{q_{1}} } + 1 \text{ cm} &= \frac{6}{x} \text{ cm} \\
x &= \tfrac{6}{ \sqrt{ \frac{q_{2}}{q_{1}} } + 1 } \text{ cm} \\
  &= \tfrac{6}{ \sqrt{ \frac{-2.0}{-8.0} } + 1 } \text{ cm} \\
  &= \frac{6}{ \frac{3}{2} } \text{ cm} = 4 \text{ cm}
\end{align*}

\newpage
\begin{align*}
\text{2. Variables: }
\text{rod length} &= d + L - d = L \\
          \lambda &= cx,\ c > 0 \\
\end{align*}
\begin{align*}
\begin{aligned}[t]
2.(a)\ \lambda &= cx \\
       \because &\ \text{ SI unit for } x \text{ is m} \\
       \therefore &\ \frac{C}{m} = c \cdot m \\
       c &= \frac{C}{m^{2}}
\end{aligned}
&\quad
\begin{aligned}[t]
2.(b)\ dQ &= \lambda \ dx \\
        Q &= \int_{d}^{d+L} \lambda \ dx \\
          &= \int_{d}^{d+L} cx \ dx \\
          &= c \int_{d}^{d+L} x \ dx \\
          &= \tfrac{1}{2} c x^{2} \bigg |_{d}^{d+L} \\
          &= \tfrac{1}{2} c \ \big [ (d+L)^{2} - d^{2} \big ] \\
          &= \tfrac{1}{2} c \ ( 2dL + L^{2} )
\end{aligned}
\end{align*}

2.(c) Since rod lays on the positive end of \(x\)-asix and the point asked for electric field is at origin, left of the the rod, therefore the unit vector points left, \(\hat{\imath} = 180^{\circ}\).
\begin{align*}
\overrightarrow{dE} &= k_{e} \frac{dq}{x^{2}} \hat{\imath} \\
\overrightarrow{E} &= k_{e} \int_{d}^{d+L} \frac{\lambda \ dx}{x^{2}} \hat{\imath} \\
                   &= k_{e} \int_{d}^{d+L} \frac{cx}{x^{2}} \ dx \hat{\imath} \\
                   &= k_{e}c \int_{d}^{d+L} \frac{x}{x^{2}} \ dx \hat{\imath} \\
                   &= k_{e}c \int_{d}^{d+L} \frac{1}{x} \ dx \hat{\imath} \\
                   &= k_{e}c \ln{x} \big |_{d}^{d+L} \hat{\imath} \\
                   &= k_{e}c \big [ \ln(d+L) - \ln(d) \big ] \hat{\imath} \\
                   &= k_{e}c \ \ln( \frac{d+L}{d}) \hat{\imath}
\end{align*}

\newpage
\begin{align*}
2.(d)\ \text{context: } \\
\text{if } L &= 2d \text{,} \\
\text{then } Q &= \frac{1}{2}c(4d^2+4d^2) \\
              &= 4 cd^2 \\
\text{and } \overrightarrow{E} &= k_e c \ (\ln \frac{d+2d}{d}) \hat{\imath} \\
\overrightarrow{E} &= k_e c \ln 3 \hat{\imath} \\
\overrightarrow{E} &= k_{e} \frac{Q}{r^{2}} \\
k_{e}c \ \ln3 \hat{\imath} &= k_{e} \frac{4cd^{2}}{r^{2}} \\
                           &= 4k_{e}c \frac{d^{2}}{r^{2}} \\
\tfrac{\ln 3}{4} \hat{\imath} &= \tfrac{d^{2}}{r^{2}} \\
\tfrac{1}{r^{2}} \hat{\imath} &= \tfrac{\ln 3}{4d^{2}} \\
                            r &= \tfrac{2d}{ \sqrt{\ln 3} } \cdot -\hat{\imath}
\end{align*}

2.(e)
For the sake of this problem and to make the unit cancellation clearer, let’s set the value of \(c\) as \(v \,\frac{C}{m^2}\).
\begin{align*}
\overrightarrow{F} &= q_{0} \overrightarrow{E} \hat{\imath} \\
                   &= p^{+} k_{e}c \ln\!\left(\frac{d+L}{d}\right) \hat{\imath} \\
                   &= 1.6 \times 10^{-19}C \cdot 8.99 \times 10^{9} 
                      \,\frac{N m^2}{C^2} \cdot v \frac{C}{m^2}
                      \ln\!\left(\frac{d+L}{d}\right)\hat{\imath} \\
                   &= 1.44 \times 10^{-9} \cdot v \ln\!\left(\frac{d+L}{d}\right) N \hat{\imath}
\end{align*}

The force by the rod onto the proton is towards the positive \(\hat{\imath}\) direction, therefore, the negative \(x\)-axis direction.
\begin{align*}
\overrightarrow{a} &= \frac{ \overrightarrow{F} }{m_{p^{+}}} \\
&= \frac{ 1.44 \times 10^{-9} \cdot v \ln\!\left(\frac{d+L}{d}\right) N }{ 1.68 \times 10^{-27} \, kg} \hat{\imath}
\end{align*}
The acceleration is not at a constant rate, with \(v \ln\!\left(\frac{d+L}{d}\right)\) being a variable.
\end{document}
